% Generated draft: requires mathematical and source review.
\documentclass[11pt]{article}
\usepackage[margin=25mm]{geometry}
\usepackage{fontspec}
\setmainfont{Latin Modern Roman}
\usepackage{amsmath,amssymb,mathtools,unicode-math}
\setmathfont{Latin Modern Math}
\usepackage{enumitem,needspace,xurl,hyperref}
\setlength{\parindent}{0pt}
\setlength{\parskip}{5pt}
\setlength{\emergencystretch}{3em}
\newcommand{\N}{\mathbb{N}}
\newcommand{\Z}{\mathbb{Z}}
\newcommand{\Q}{\mathbb{Q}}
\newcommand{\R}{\mathbb{R}}
\newcommand{\C}{\mathbb{C}}
\newcommand{\F}{\mathbb{F}}
\newcommand{\A}{\mathbb{A}}
\newcommand{\PP}{\mathbb{P}}
\newcommand{\E}{\mathbb{E}}
\newcommand{\eps}{\varepsilon}
\newcommand{\dd}{\,\mathrm{d}}
\begin{document}
\section*{UnsolvedMath ID: 20002037}
\textbf{Source problem code:} catalogue entry
\subsection{Definitions and mathematical statement}
Fix a connected oriented time-oriented Lorentzian conformal spin manifold \((M^n,[g])\), \(n\geq3\), with a chosen spin structure. Let \(h(g)=\dim_\C\ker D_g\) and \(t(g)=\dim_\C\ker P_g\), for global smooth Dirac and twistor spinors respectively. Problem 17 asks to describe the conformally flat structures for which \(h(g)>0\) or \(t(g)>0\). The missing definitions are supplied by companion record 20002032; no \(L^2\) or support restriction is imposed in this version.

For the twistor branch the supplied aid reports a full-monodromy reduction. On a conformally flat manifold the conformal spin-tractor connection is flat and twistor spinors correspond to its parallel sections. If \(V\) is its complex spin-tractor fiber and \(\rho:\pi_1(M)\to\operatorname{GL}(V)\) its full lifted monodromy representation, then
\[
\ker P_g\cong V^{\rho(\pi_1(M))},\qquad
t(g)=\dim_\C V^{\rho(\pi_1(M))}.
\]
The spin lift is part of the data. Restricted holonomy is trivial for this flat connection and alone does not decide global descent.

For a flat Lorentzian torus \(\R^n/\Lambda\), let \(\chi:\Lambda\to\{\pm1\}\) specify the spin structure, choose \(\alpha\in\tfrac12\Lambda^*/\Lambda^*\) with \(\chi(\lambda)=e^{2\pi i\alpha(\lambda)}\), and use frequencies \(k\in\Lambda^*+\alpha\). The aid gives the Fourier description
\[
\psi(x)=\sum_k a_ke^{2\pi i k(x)},\qquad
k^{\sharp_g}\cdot a_k=0,
\]
with coefficients rapidly decreasing in any auxiliary Euclidean frequency norm. Only frequencies on the metric null cone contribute. At \(k=0\) the coefficient is unrestricted; at nonzero null \(k\) the Clifford kernel has dimension \(N/2\), where \(N=2^{\lfloor n/2\rfloor}\). Thus the harmonic kernel is determined by the shifted dual lattice and the null cone, together with this rapid-decay condition. The reported twistor dimensions are \(t=N\) for trivial \(\chi\) and \(t=0\) for nontrivial \(\chi\).

The aid's four dimension pairs are realized concretely as follows. On the unit-period torus with \(g=-\dd t^2+\sum_i\dd x_i^2\), trivial spin structure gives \((h,t)=(\infty,N)\); antiperiodicity only in time gives \((0,0)\); antiperiodicity in time and one spatial coordinate gives \((\infty,0)\), since corresponding half-integral null frequencies exist. With trivial spin structure and \(g=-\dd t^2/\sqrt2+\sum_i\dd x_i^2\), the dual null equation \(\sum_i k_i^2=\sqrt2 k_0^2\) has no nonzero integral solution, giving \((N,N)\). These examples show that \(h\) and \(t\) have different global analytic behavior.

The aid separately reports that on globally hyperbolic Lorentzian spin manifolds the unrestricted smooth Dirac kernel is infinite dimensional by the Cauchy problem. Compact Lorentzian tori are not examples of that globally hyperbolic setting. The monodromy and Fourier descriptions are supplied reductions and a classified family, not a complete explicit classification of all conformally flat Lorentzian conformal spin manifolds.

The extraction's following section, from a \(Q\)-curvature primer, is unrelated spillover and is not used as a hypothesis or reconstructed as part of Problem 17.
\subsection{Short English statement}
Describe conformally flat Lorentzian spin manifolds with harmonic or twistor spinors. The supplied aid reduces twistor existence to full lifted monodromy and computes flat-torus kernels from spin-shifted null frequencies. Spin structure and the global smooth function space remain essential.
\subsection{Sources}
\begin{itemize}
\item Helga Baum, Problem 17, AIM workshop \emph{Conformal structure in geometry, analysis, and physics}: \url{https://aimath.org/WWN/confstruct/confstruct.pdf}.
\item Local registry \texttt{data/unsolvedmath/problems.json}, permanent ID 20002037, code \texttt{catalogue entry}; canonical \texttt{aim-geometry-notes.json}, \texttt{notes[374]}; updated 14 August 2026. Companion ID 20002032, code \texttt{AIM-GEOMETRY-0370}, \texttt{notes[369]}, supplies Baum's definitions of \(h,t\).
\item The record's unrecoverable statement-status label concerns merged extraction. Its verification separates Problem 17 from the next section, \emph{A Primer on Q-curvature} by M. Eastwood and J. Slov\'ak; that material is not part of this problem.
\item Attempt 1: partial result, problem status at run open, rigor label proved, novelty confidence low; current status partially solved. No full external reference for the spin-tractor or Fourier claims is supplied; the original literature field is empty. The aid is machine generated, not peer reviewed, and requires independent verification before citation.
\end{itemize}
\end{document}
